\scanpage{129}
\deleted{$\unsure{g_i}$}

\deleted{\unsure{(8)}} Si le graphe est MD, alors $(g,\nu)$ est admissible, et on pose
\[
(9)\qquad \unsure{\delta|d}(G)=\unsure{3}g-3+\nu
\qquad (\uline{\text{dimension}}\ \text{modulaire}\inserted{\uline{ambiante}})
\]
et on aura
\[
(10)\qquad d(G)+c(G)=\delta(G)
\]

\uline{Opérations} \inserted{\uline{et morphismes}.}
Soit $G,G'$ \deleted{deux} graphes MD, on appelle \uline{hom de générisation} \inserted{$u$} de $G$ dans $G'$ un couple d'applications
\diagram{\begin{tikzcd}[row sep=small]
S\arrow[r,"u_s"]&S'\\
\widehat{\vec A}\arrow[r,"u_a"]&\widehat{\vec A}'\amalg S'
\end{tikzcd}}
satisfaisant les conditions suivantes, où
\[
\vec B=u_a^{-1}(\widehat{\vec A}'),\qquad
\vec C=\widehat{\vec A}\setminus\unsure{B}=u_a^{-1}(S')
\]
a) $\vec B\supset I$ (ens des arcs libres), $\vec B$ stable par $\sigma$

b) L'application $\vec B\longrightarrow\widehat{\vec A}'$ induite par $u_a$ commute aux involutions et le couple $(u_s,u_a|\vec B)$ \unsure{commutent|commute} aux applications origine.
(Donc $u_a$ induit $I\longrightarrow I'$)

\margin{i.e. $(S,\vec B)$ \deleted{\unsure{a}} définit un sous-graphe de $(S,\widehat{\vec A},o,\sigma)$ contenant les arêtes libres de $G$; $(u_s,u_a|\vec B)$ définit un hom. de graphes.}

c) \deleted{$\unsure{a\in\vec C}$} si $c\in\vec C$, alors
\[
u_a(c)=\inserted{u_a(-c)=}u_s(\underbrace{o(c)}_{\text{origine}})
\]
(donc $u_a|C$ est connu quand on \unsure{connaît|connait} $u_s$)\deleted{\unsure{a}}. \unsure{donc}
[Donc \deleted{\unsure{de $u_s$ doit}} $u_s$ a même image sur deux \deleted{\unsure{les}} sommets $s,s'$ tels que $\exists c\in\vec C$, $s=o(c)$, $s'=o(-c)$ --- donc $u_s$ se factorise en
\diagram{\begin{tikzcd}
S\arrow[r]&S/R\arrow[r]&S'
\end{tikzcd}}
où $R$ est la relation d'équivalence dans $S$ engendrée par la relation précédente]

d) $S/R\longrightarrow S'$ induit par $u$ est une bijection,
$\vec B\longrightarrow\widehat{\vec A}'$ est une bijection

\margin{donc $S'$ et $\widehat{\vec A}'$ se récupèrent à partir de $\vec B\subset\widehat{\vec A}$ i.e. de $C$, à bijection can. près\ldots}

e) Pour tout $s'\in S'$, soit $S(s')=u_s^{-1}(\{s'\})$, $\unsure{\vec C(s')|\widehat{\vec A}(s')}=$ ens. des $\alpha\in\deleted{\unsure{\widehat A}}\vec C$ tels que
\deleted{$u_a(\alpha)=s'\ \text{ou}\ \unsure{u_a(\alpha)}$}
$o(\alpha)$ et $e(\alpha)\in S(s')$
(i.e. $u_a(\alpha)=s'$)
\deleted{ou $\alpha\in B$ et $\unsure{o'(u_a(\alpha))=e'(u_a(\alpha))=s'}$)}.
Soit \deleted{alors} $G_0(s')$ le sous-\deleted{\unsure{MD}}-graphe \inserted{pondéré} de $G$ défini par $S(s')$ et $\unsure{\widehat{\vec A}(s')}$, considérons son genre virtuel
\[
\uline{g}(G_0(s'))=\sum_{s\in S(s')}g_s+h^1(G_0(s'))
\]
\margin{\uline{NB} d) assure que $G(s')$ est connexe}
On a alors
\[
\uline{g}_{s'}=g(G_0(s'))
\]
\mspage{128}
